\begin{lemma}[Uniqueness of the weak gradient in $L^2(\gamma_n)$]
Let $v : \mathbb{R}^n \to \mathbb{R}$ be measurable and let
$G, H \in L^2(\gamma_n;\mathbb{R}^n)$ both be weak Euclidean gradients of $v$ in
the sense of \noderef{HasGaussWeakGrad}.  Then $G = H$ as elements of
$L^2(\gamma_n;\mathbb{R}^n)$, i.e. $G = H$ $\gamma_n$-almost everywhere.
\end{lemma}
\begin{proof}
Fix $j \in \{1,\ldots,n\}$ and put $w = G_j - H_j$.

\emph{Step 1: $w$ annihilates every test function.}
By \noderef{HasGaussWeakGrad} applied to the pair $(v,G)$ and then to the pair
$(v,H)$, for every $\varphi \in C_c^\infty(\mathbb{R}^n)$,
\[
  \int v\,\partial_j\varphi\, dx = -\int G_j\,\varphi\, dx,
  \qquad
  \int v\,\partial_j\varphi\, dx = -\int H_j\,\varphi\, dx ,
\]
the integrals being with respect to Lebesgue measure.  Subtracting,
\begin{equation*}
  \int_{\mathbb{R}^n} w\,\varphi\, dx = 0
  \qquad\text{for every } \varphi \in C_c^\infty(\mathbb{R}^n). \tag{1}
\end{equation*}

\emph{Step 2: $w$ is locally Lebesgue integrable.}
Both $G_j$ and $H_j$ lie in $L^2(\gamma_n)$, since $G$ and $H$ lie in
$L^2(\gamma_n;\mathbb{R}^n)$.  By \noderef{GaussL2LocallyIntegrable}, applied to
each of them, each is locally integrable for Lebesgue measure.  A difference of
two locally Lebesgue-integrable functions is locally Lebesgue integrable, so
$w = G_j - H_j \in L^1_{\mathrm{loc}}(\mathbb{R}^n)$ for Lebesgue measure.

\emph{Step 3: $w = 0$ Lebesgue-a.e.}
The du Bois-Reymond item of shared context says exactly that a function which is
locally integrable for Lebesgue measure and integrates to zero against every
$\varphi \in C_c^\infty(\mathbb{R}^n)$ vanishes Lebesgue-almost everywhere.  Its
two hypotheses are in hand: local Lebesgue integrability of $w$ is Step 2, and
the vanishing of $\int w \varphi \, dx$ for every such $\varphi$ is $(1)$ from
Step 1.  Hence $w = 0$ Lebesgue-almost everywhere.

\emph{Step 4: transfer to $\gamma_n$.}
Again by \noderef{GaussMeasure}, $\gamma_n$ has a strictly positive density with
respect to Lebesgue measure, so a Lebesgue-null set is $\gamma_n$-null and
conversely.  Therefore $G_j = H_j$ $\gamma_n$-almost everywhere.

As $j$ was arbitrary, $G = H$ $\gamma_n$-almost everywhere, i.e.\ $G = H$ in
$L^2(\gamma_n;\mathbb{R}^n)$.
\end{proof}
